\phantomsection
\addcontentsline{toc}{section}{Bảng thuật ngữ}
\section*{Bảng thuật ngữ}

\begin{table}[htbp]
    \centering
    \caption{Bảng đối chiếu thuật ngữ Anh - Việt}
    \label{tab:bang-thuat-ngu}
    \begin{tabularx}{\textwidth}{l X}
        \toprule
        \textbf{Thuật ngữ tiếng Anh}       & \textbf{Giải nghĩa tiếng Việt}                                                                \\
        \midrule
        Attention                          & Cơ chế tập trung (cơ chế chú ý giúp mô hình tập trung vào các vùng thông tin quan trọng)      \\
        Annotation                         & Gán nhãn dữ liệu (chú thích thông tin lớp và vị trí đối tượng trên ảnh)                       \\
        Batch size                         & Kích thước lô (số lượng mẫu dữ liệu được xử lý trong một lần cập nhật trọng số)               \\
        Bounding box                       & Khung bao (hộp giới hạn xác định vị trí đối tượng trong ảnh)                                  \\
        Convolutional Neural Network (CNN) & Mạng thần kinh tích chập (học đặc trưng hình ảnh cục bộ thông qua phép tích chập)             \\
        Data Augmentation                  & Tăng cường dữ liệu (tạo thêm dữ liệu huấn luyện bằng các phép biến đổi hình thái, màu sắc)    \\
        Dataset                            & Tập dữ liệu (nguồn dữ liệu sử dụng để huấn luyện, kiểm tra và đánh giá mô hình)               \\
        Edge devices                       & Thiết bị biên (thiết bị phần cứng đầu cuối có tài nguyên tính toán giới hạn)                  \\
        Epoch                              & Vòng lặp huấn luyện (một chu kỳ mô hình học qua toàn bộ tập dữ liệu huấn luyện)               \\
        Gradient accumulation              & Tích lũy đạo hàm (kỹ thuật gom đạo hàm qua nhiều bước nhỏ để mô phỏng kích thước lô lớn)      \\
        Human-in-the-loop                  & Vòng lặp có sự tham gia của con người (AI đề xuất kết quả, con người kiểm duyệt và chỉnh sửa) \\
        Hyperparameter tuning              & Tinh chỉnh siêu tham số (tối ưu hóa các tham số thiết lập trước khi huấn luyện)               \\
        Instance Segmentation              & Phân vùng thực thể (phân lớp từng điểm ảnh và phân biệt rõ các thực thể độc lập)              \\
        Learning rate                      & Tốc độ học (bước nhảy cập nhật trọng số trong thuật toán tối ưu)                              \\
        Letterbox                          & Chuẩn hóa kích thước ảnh giữ nguyên tỷ lệ khung hình thực (thêm viền đệm)                     \\
        Loss function                      & Hàm mất mát (hàm đánh giá độ lệch giữa dự đoán của mô hình và nhãn thực tế)                   \\
        Mask                               & Mặt nạ (vùng điểm ảnh nhị phân xác định hình thái chính xác của thực thể)                     \\
        Overfitting / Underfitting         & Quá khớp / Chưa khớp (tình trạng ghi nhớ quá mức hoặc không đủ khả năng học đặc trưng)        \\
        Prompt                             & Gợi ý đầu vào (thông tin dẫn hướng như điểm, khung bao giúp mô hình phân vùng)                \\
        Stratified Sampling                & Lấy mẫu phân tầng (kỹ thuật chia tập dữ liệu đảm bảo đồng đều tỷ lệ các lớp)                  \\
        Vision Transformer (ViT)           & Kiến trúc Transformer cho thị giác máy tính (học đặc trưng toàn cục nhờ cơ chế Attention)     \\
        Warmup                             & Khởi động ấm (tăng dần tốc độ học từ nhỏ đến giá trị mục tiêu trong các epoch đầu)            \\
        Weight decay                       & Suy giảm trọng số (kỹ thuật chính quy hóa giúp ngăn chặn tình trạng quá khớp)                 \\
        Zero-shot                          & Khả năng suy luận trực tiếp trên dữ liệu mới mà không cần huấn luyện bổ sung                  \\
        \bottomrule
    \end{tabularx}
Dưới đây là danh sách các thuật ngữ viết tắt và các thuật ngữ chuyên ngành tiếng Anh được sử dụng trong báo cáo đồ án, kèm theo thuật ngữ dịch nghĩa hoặc giải nghĩa trong tiếng Việt.

\begin{table}[htbp]
\centering
\caption{Bảng đối chiếu thuật ngữ Anh - Việt}
\label{tab:bang-thuat-ngu}
\begin{tabularx}{\textwidth}{l X}
\toprule
\textbf{Thuật ngữ tiếng Anh} & \textbf{Giải nghĩa tiếng Việt} \\
\midrule
Annotation & Gán nhãn hoặc chú giải dữ liệu (tọa độ, nhãn lớp hoặc mặt nạ phân vùng). \\
Aspect ratio & Tỉ lệ khung hình (tỉ lệ giữa chiều rộng và chiều cao của ảnh). \\
Bounding box & Hộp giới hạn bao quanh đối tượng hoặc vùng bệnh được phát hiện. \\
Classification & Bài toán phân loại ảnh (gán một hoặc nhiều nhãn cho toàn bộ bức ảnh). \\
Data augmentation & Tăng cường dữ liệu (các phép biến đổi hình học hoặc quang học ảnh). \\
Dataset & Bộ dữ liệu (bao gồm tập ảnh và nhãn đi kèm). \\
Domain shift & Sự dịch chuyển miền dữ liệu (khác biệt phân phối giữa dữ liệu học và thực tế). \\
Human-in-the-loop & Quy trình có con người tham gia kiểm duyệt và tinh chỉnh kết quả của AI. \\
Instance Segmentation & Phân vùng thực thể (tách biệt từng đối tượng ở cấp độ pixel). \\
mAP (Mean Average Precision) & Độ chính xác trung bình trung (chỉ số đánh giá phát hiện vật thể). \\
mIoU (Mean Intersection over Union) & Tỷ lệ giao trên hợp trung bình (chỉ số đánh giá độ trùng khớp mặt nạ). \\
ONNX (Open Neural Network Exchange) & Chuẩn trao đổi mạng nơ-ron mở giúp tối ưu hóa suy luận mô hình. \\
VLM (Vision Language Model) & Mô hình thị giác - ngôn ngữ (mô hình đa phương thức xử lý ảnh và chữ). \\
ASHA (Asynchronous Successive Halving) & Thuật toán dừng sớm bất đồng bộ phục vụ tìm kiếm siêu tham số tối ưu. \\
Dice Score & Chỉ số đo lường độ trùng khớp giữa hai vùng mặt nạ (F1-score ở cấp pixel). \\
WeightedRandomSampler & Bộ lấy mẫu ngẫu nhiên có trọng số dùng để xử lý mất cân bằng lớp. \\
\bottomrule
\end{tabularx}
\end{table}
